\documentclass{article}
\usepackage{amsmath,amssymb,hyperref}
\title{An exclusive shared brake can make a formation certificate false}
\author{}
\date{26 November 2026}
\begin{document}
\maketitle
\paragraph{Pins.}
After a fault, two vehicles traveling along a straight corridor must stop
at specified checkpoints without losing clearance or geometric radio reach.
They still have one shared braking-capacity unit. One design can split it;
another awards the entire unit to one vehicle for each fixed-duration
slot, or idles. This is an actuator model, not data about a deployed
platform or an RF reception claim. Positions and velocities are known;
the obstruction is physical in the stated control class, not missing
evidence. The fault does not itself violate the current gap; only the
available braking interface changes. Constant nonnegative deceleration,
no acceleration or reverse travel, no disturbance after the fault, and
fixed slot boundaries are stipulated.

In units with slot duration and acceleration limit both one, the two
initial speeds are $v>0$ and the deadline is $T=2v$. Let $q$ denote
front-minus-rear gap change from the fault cut. Clearance is $q\ge-D$
and geometric communication requires $q\le U$ throughout, where
both slacks are nonnegative relative to the initial gap. The terminal
gap change $k$ fixes two individual distances when both stop:
their sum is $2v^2$, so the front and rear go $v^2+k/2$ and $v^2-k/2$.
For example $k=0$ meets both individual caps $v^2$ exactly. Any
fixed-obstacle cap on these monotonically traveled distances can
therefore be checked directly. A relaxed post-fault gap is a changed
requirement, not evidence of having preserved the original one.

\paragraph{Exact split-capacity comparison.}
The deadline and nonnegative velocities force all capacity to be spent
almost everywhere: total velocity is $2v-t$. Thus total travel is
$2v^2$. Each individual speed is between $\max(0,v-t)$ and
$\min(v,2v-t)$ by its deceleration bound and its need to stop by $T$.
Hence $|k|\le v^2$. The extreme schedules award the first $v$ units
of time solely to one vehicle, then the next $v$ to the other. Their
relative-velocity deviations have a consistent sign until stop.
Convex combinations of the two schedules remain physically admissible
under splittable actuation and make each intermediate $k$ with a monotone
gap change. Thus the split interface meets the entire tube precisely
when $|k|\le v^2$ and $k\in[-D,U]$. In particular an equal split
keeps $q=0$ precisely, not only at the sampled stops.

\paragraph{Exact exclusive-slot law near zero terminal displacement.}
Now require $v\in\mathbb N$ and award the sole channel to one vehicle
on each unit slot. Neither idling nor braking a stopped vehicle meets
the deadline; exactly $v$ of the $2v$ slots go to each vehicle. For
integer targets $|k|\le v$, a trajectory exists if and only if
\[
 k\in[-D,U],\qquad k\equiv v\pmod 2,\qquad \max(D,U)\ge1 .
\]
The last condition is necessary independently of the target.
The continuous gap at time two is $(3a_0+a_1)/2$, where
$a_i\in\{+1,-1\}$ is the signed gap acceleration. It has magnitude
at least one. For parity, integrating all slots gives
$k=\sum_{j=0}^{2v-1}(2v-j-1/2)a_j$. The signed halves cancel.
The residual integer is twice the sum of weights with positive sign
minus $v(2v+1)$, hence has the asserted parity.

For sufficiency use two-slot blocks with opposite assignments.
Each returns relative velocity to zero, brakes both vehicles once,
and moves the gap monotonically by one unit in a chosen direction.
Use $|k|$ such moves toward $k$ and $(v-|k|)/2$ excursions that return
to zero, inserting the latter first. If $k\ne0$, the endpoint condition
allows a one-unit excursion on its side. If $k=0$, choose either side
with slack at least one. The construction stays inside the tube at every
continuous time. The claim does not characterize arbitrary $|k|>v$,
unequal initial speeds, or changeable slot lengths during a run.

\paragraph{Consequence and baselines.}
At $v=2$, $D=U=1/2$, the terminal constraint alone admits exclusive
schedules but none keeps the tube. Giving the increasing-gap side
one unit suffices without weakening the clearance lower bound. At
odd $v$, exact restoration of $k=0$ is impossible for the exclusive
interface, even with a very wide tube. In physical units, the excursion
floor is $a h^2$ for deceleration $a$ and slot duration $h$; the parity
test applies to $v=V/(a h)$ for physical speed $V$ in this family.
With $V=2,a=1$, halving the hold from one to one-half preserves the
deadline and both checkpoint distances while reducing the witnessed
excursion from one to one-quarter. The original two half-unit slacks
then suffice without changing clearance or radio geometry.

\href{https://arxiv.org/html/2210.01354v2}{Breeden and Panagou,
Section II-B} assume a compact convex admissible input set for
their joint-CBF method. A certificate after convexifying the
slot-exclusive set proves a statement about a different actuator,
not a defect in that result. The convex comparison here is exact.
\href{https://arxiv.org/html/2609.29939v1}{Safe Formation Control of
Open Multi-Robot Systems} treats managed prospective edges before
departure; no such warning or prospective physical capacity is added
in this post-fault comparison. Maeda--Endo--Matsuno's decentralized
crash-recovery comparison remains open in this study.
Exact direct enumeration agrees on 2268 band/target queries and
checks 4706 exclusive schedules with within-slot integration. Four
Lean statements check the key algebraic kernels. This isolates a
controller-to-actuator boundary in R6, not full R6 acceptance or
any independent mission result.
\end{document}
